\section{Interpretation boundaries}\label{sec:limits}
The following distinctions are part of the mathematical result, not qualifications to be removed by a more suggestive picture.
\begin{center}\small
\begin{tabularx}{\linewidth}{@{}XcX@{}}\toprule
Established mathematical object & & Unsupported identification\\\midrule
Response inner product &\(\ne\)& Spacetime metric\\
Screened graph response &\(\ne\)& Newtonian gravity\\
Complex response eigenmode &\(\ne\)& Gravitational wave\\
Non-reversible cycle imbalance &\(\ne\)& Physical current\\
SPD symmetrizer &\(\ne\)& Privileged physical geometry\\
Tangent vector &\(\ne\)& Physical position\\
Spatial observable &\(\ne\)& Core trajectory\\
Axial parity &\(\ne\)& Physical rotation axis\\
State-to-face representation &\(\ne\)& Material deformation\\
Graph coupling &\(\ne\)& Seam flux\\
Lens/SRG preparation &\(\ne\)& Boundary dynamics\\\bottomrule
\end{tabularx}
\end{center}
The inequality signs denote unestablished identifications, not an assertion that no future model could relate the objects. The response inner product depends on a state background; the panel metric and matched connection are fixed. Their different domains prevent an automatic identification. Likewise, the non-real eigenvalues are statements about powers of a finite Jacobian, without an assigned physical clock or wavelength.

For nonnegative response entries, a row-sum-one matrix has a mathematical finite-state Markov interpretation, and the cycle obstruction tests reversibility. Without entrywise nonnegativity the algebraic balance question still makes sense, but probability currents should not be asserted. Even a valid Markov interpretation would not itself establish a physical current.

The names of the GR checkpoints preserve the original investigation question. Their conclusion is response geometry and finite spectral structure, not a physical gravity theory. No new physical model, kernel integration, historical-motion import or parked-lane derivation is part of this paper.

\section{Relation to established mathematics}\label{sec:discussion}
Equation~\eqref{eq:weighted} belongs to the weighted graph Dirichlet-form setting described by Chung \cite[Chapter~1]{Chung}. The reversible/non-reversible distinction is an edge-balance problem; in a positive-entry stochastic domain, the triangle proof is the finite Kolmogorov cycle condition. Detailed balance and reversibility are standard Markov-chain concepts \cite[\S1.6]{LP}; the matrix cycle condition is also treated in \cite[\S4]{MS}. The native cycle factor and its exceptional parameter loci are derived here from the checkpoint equations, not imported from these references.

The terminology in Section~\ref{sec:spd} separates this node-local condition from arbitrary SPD self-adjointness. McKee and Smyth \cite[Definition~4.1]{MS} use ``symmetrizable'' for similarity to a symmetric matrix by a positive diagonal change of basis. Theorem~\ref{thm:spd}, proved here, permits a general SPD inner product and is correspondingly weaker than diagonal balance.

The exact translation-by-three decomposition is finite Fourier or Bloch analysis. Kuchment's overview \cite[\S4]{Kuchment} describes Floquet decomposition in the broader theory of periodic operators, including operators on graphs. Our finite blocks are derived directly; that context supplies neither a continuum approximation nor a physical wavelength for the native response.

The matched maps identify specified vector fibres compatibly. Graph connection and vector-fibre terminology has an established setting in Singer and Wu \cite{SW}. This comparison does not identify the native transport with their data-derived connection Laplacian or infer an ambient gravitational connection. Likewise, the observable census uses finite-dimensional equivariant polynomial representation theory; Golubitsky, Stewart and Schaeffer \cite{GSS} provide general group-theoretic and equivariant-map context. The native representation, multiplicities and observable formulas remain derived in this paper and its cited checkpoints.

The screened response is a finite modified-Helmholtz inverse. The closure theorem is the descent-to-a-quotient criterion, proved here for the native recurrence with its zero convention. These comparisons assign mathematical types; they assert neither novelty nor a physical dictionary. The internal accepted bibliography supports the native definitions and prior ownership \cite{PA,PB,PC,PD,PE,PF,AT04,AT08}. The six external references supply targeted surrounding context only. All model-specific claims and the linear-algebra, Fourier, cycle and quotient arguments needed here remain self-contained; no priority claim is made.

\section{Limitations, open questions and conclusion}
Every nonlinear branch existence and attraction statement here is local for a chosen finite ring unless explicitly stated otherwise. The positive-branch algebra can have a broader domain than its local stability consequences. The generic six-ring neighbourhood is not uniform near the nilpotent fixture. The cubic imaginary-root expansion requires a simple nonzero baseline root; exact polynomial and rank tests cover the exceptional cases.

The finite-carrier census fixes the full \(D_{3h}\) representation and polynomial degree. It does not classify arbitrary nonpolynomial maps, all continuous shell fields or observations depending jointly on evolving additional geometry. The coherence quotient is regular-domain mathematics, with the native zero-stratum counterexample retained. No proof of a complete nonlinear continuum limit, global positive-branch continuation, canonical dense response metric or physical position selection is supplied.

One bounded open problem already identified by SA0 is the minimal subset of existing pairwise coherences that, together with face intensities, supports autonomous regular-domain closure. Intensity is insufficient; full real Gram is sufficient. This synthesis does not begin that calculation.

The established chain is nevertheless substantial. Background inequality selects exact pre-stage weights and conductances; synchronization generically removes diagonal reversibility. The spectral consequence depends on the finite sector: local positive triads permit a non-diagonal SPD metric; for \(N\ge9\), three-distinct repeated backgrounds satisfying the stated nondegeneracy condition force non-real modes; and a six-ring admits a native square-zero defect with a sign-resolved unfolding. Independently, the spatial representation has multiple equivariants, while its adopted complete face view is lossless. The coherence quotient explains a specific level at which passive information can evolve autonomously and identifies exactly why the native zero convention blocks the unrestricted statement. These results describe the existing mathematical map without a new physical premise.

\section{Reproducibility, evidence and availability}\label{sec:repro}
The source checkpoint is
\[
\texttt{da461294d724997bd7033c7cea9cc9a0b5f5031d}.
\]
The scientific source checkout, existing papers, Atlas and all fifteen GR0/GR1/GR2/CM0/SA0 artifacts are preserved. Only this new publication lane is authored. Source hashes, the initial dirty status, final preservation comparison, theorem ledger and paper-source crosswalk are included in the package. Private machine locators occur only in provenance records, not in scientific statements.

Retained counts are
\begin{center}
\begin{tabular}{lrrrrrr}\toprule
Checkpoint & GR0 & GR1 & GR2 & CM0 & SA0 & total\\\midrule
Exact/source checks &45&38&42&78&226&429\\\bottomrule
\end{tabular}
\end{center}
These heterogeneous predicates include symbolic identities, matrix-entry groups, source-dependency checks and bounded structural witnesses. They are not 429 theorems. Written arguments are the proof authority. Predecessor scripts are referenced, not rerun into their result files; their verification records are retained with unchanged source hashes.

Numerical evidence is kept by method and quantity. Native complex128 finite-difference errors are not combined with 100-digit transcription errors, branch residuals, characteristic residuals, metric balance residuals or spectral imaginary parts. In particular, a native Jacobian residual does not resolve imaginary roots smaller than that residual. Appendix~\ref{app:verification} summarizes these separate measurements; the complete unaggregated records are in the companion retained-evidence JSON.

Paper-local verification independently checks displayed factorization, polynomial, finite-block, nilpotent and quotient identities and validates manuscript/source references. Its count and results are separate from 429. Figure sources generate five vector schematics with explicit derivation notes. Build commands, a proposed file allowlist, publication text and three optional public-post formulations are supplied; none is published automatically. Research publication v0.1 carries no claim of journal submission, external peer review, a repository release or DOI registration.

\paragraph{Preparation disclosure.}
The manuscript was assembled with AI-assisted mathematical writing, source reconciliation, symbolic verification and typesetting. Scientific review, author approval and any publication decision remain with the project owner. No claim of independent peer review is made.
